\lecture{6.2}{Manual Cross-Validation}
\label{lec:6-2}

\subsection{Full-Data Risk vs Cross-Validated Risk}

Before calling \texttt{SuperLearner()}, it is useful to implement cross-validation by
hand. This makes clear what the package automates.

Consider two candidate estimators for the outcome regression $\Qbar(A,W)$:
\[
\Qbar_{n,1}: \text{ main-terms linear regression},
\]
\[
\Qbar_{n,2}: \text{ a more flexible regression with quadratic and interaction terms}.
\]
The mean squared-error risk of a fitted estimator $Q$ is
\[
R(Q)=\E\left[\{Y-Q(A,W)\}^2\right].
\]
The naive full-data empirical risk is
\[
R_n(Q_n)=\frac{1}{n}\sum_{i=1}^n\{Y_i-Q_n(A_i,W_i)\}^2.
\]
This quantity can be too optimistic because the same observations are used both to fit
the regression and to evaluate its error.

\begin{implementationbox}
\begin{lstlisting}[language=R]
Q1_fit <- glm(Y ~ W1 + W2 + A, data = full_data)
Q1n <- as.numeric(predict(Q1_fit))
risk_Q1_full <- mean((Y - Q1n)^2)

Q2_fit <- glm(Y ~ I(W1^2) + I(W2^2) + W1 * W2 + W1 * A + W2 * A,
              data = full_data)
Q2n <- as.numeric(predict(Q2_fit))
risk_Q2_full <- mean((Y - Q2n)^2)
\end{lstlisting}
\end{implementationbox}

\subsection{Two-Fold Splitting}

Two-fold cross-validation splits the data into two validation sets. Each fold is used
once for validation and once as part of the training data. With two folds, the training
set for fold 1 is fold 2, and the training set for fold 2 is fold 1.

\begin{implementationbox}
\begin{lstlisting}[language=R]
split <- c(rep(1, n / 2), rep(2, n / 2))

split1 <- subset(full_data, split == 1)
split2 <- subset(full_data, split == 2)
\end{lstlisting}
\end{implementationbox}

In real analyses, folds are usually randomized rather than assigning the first half to
one fold and the second half to another. For teaching, a deterministic split makes the
logic easier to see.

\subsection{Training Regressions}

For $\Qbar_{n,1}$, fit the main-terms model on split 1 and validate on split 2:

\begin{implementationbox}
\begin{lstlisting}[language=R]
Q1_fit_split1 <- glm(Y ~ W1 + W2 + A, data = split1)
Q1n_split2 <- predict(Q1_fit_split1, newdata = split2)
risk_Q1_split2 <- mean((split2$Y - Q1n_split2)^2)
\end{lstlisting}
\end{implementationbox}

Then reverse the roles:

\begin{implementationbox}
\begin{lstlisting}[language=R]
Q1_fit_split2 <- glm(Y ~ W1 + W2 + A, data = split2)
Q1n_split1 <- predict(Q1_fit_split2, newdata = split1)
risk_Q1_split1 <- mean((split1$Y - Q1n_split1)^2)
\end{lstlisting}
\end{implementationbox}

\subsection{Validation Predictions}

The essential rule is that validation predictions must be made using a model that did
not train on the validation observations. For observation $i$ in validation fold $v$,
the prediction has the form
\[
\Qbar_{n,v}(A_i,W_i),
\]
where $\Qbar_{n,v}$ was fit using the training observations outside fold $v$.

\begin{figure}[H]
\centering
\begin{tikzpicture}[
  node distance=1.35cm,
  box/.style={draw, rounded corners=2pt, minimum width=3.0cm, minimum height=.75cm},
  arr/.style={-{Latex[length=2.5mm]}, thick}
]
\node[box] (s1) {Split 1};
\node[box, right=of s1] (fit1) {Fit model};
\node[box, right=of fit1] (pred2) {Predict\\ Split 2};
\draw[arr] (s1) -- (fit1);
\draw[arr] (fit1) -- (pred2);
\node[box, below=1cm of s1] (s2) {Split 2};
\node[box, right=of s2] (fit2) {Fit model};
\node[box, right=of fit2] (pred1) {Predict\\ Split 1};
\draw[arr] (s2) -- (fit2);
\draw[arr] (fit2) -- (pred1);
\end{tikzpicture}
\caption{Two-fold cross-validation swaps training and validation roles.}
\label{fig:ch6-twofold}
\end{figure}

\subsection{Mean Squared-Error Risk}

The two-fold cross-validated risk estimate is the average of the two fold-specific
validation risks:
\[
R_{CV,n}(\Qbar_{n,1})
=
\frac{1}{2}\left\{
R_{n,1}(\Qbar_{n,1,2})+
R_{n,2}(\Qbar_{n,1,1})
\right\}.
\]
In code:

\begin{implementationbox}
\begin{lstlisting}[language=R]
cv_risk_Q1 <- (risk_Q1_split1 + risk_Q1_split2) / 2
\end{lstlisting}
\end{implementationbox}

For the flexible model, repeat the same process:

\begin{implementationbox}
\begin{lstlisting}[language=R]
Q2_formula <- Y ~ I(W1^2) + I(W2^2) + W1 * W2 + W1 * A + W2 * A

Q2_fit_split1 <- glm(Q2_formula, data = split1)
Q2n_split2 <- predict(Q2_fit_split1, newdata = split2)
risk_Q2_split2 <- mean((split2$Y - Q2n_split2)^2)

Q2_fit_split2 <- glm(Q2_formula, data = split2)
Q2n_split1 <- predict(Q2_fit_split2, newdata = split1)
risk_Q2_split1 <- mean((split1$Y - Q2n_split1)^2)

cv_risk_Q2 <- (risk_Q2_split1 + risk_Q2_split2) / 2
\end{lstlisting}
\end{implementationbox}

\subsection{Comparing Simple and Flexible Regressions}

The candidate with lower cross-validated risk has better estimated out-of-sample
prediction performance under the chosen loss. A flexible model can have lower
training risk but higher cross-validated risk if it overfits. Conversely, a simple model
can win when the true signal is simple or when the sample size is limited.

\begin{keyidea}
Cross-validation evaluates how well a regression predicts observations it did not use
for training. This is why cross-validated risk is a better guide for selecting nuisance
regressions than full-data training error.
\end{keyidea}
